% ==============================================================================
% CHAPITRE 4 -- RESTITUTION DES RÉSULTATS, CONFRONTATION ET RECOMMANDATIONS
% Mémoire MAE -- Arthur SAUNIER -- Wavestone Luxembourg
% Le \chapter{} est déclaré dans le fichier maître ; ce fichier démarre en \section.
% ==============================================================================

\section{Restitution du matériau}

La restitution suit les axes de la grille reproduite en annexe~B et non l'ordre des
répondants. La matrice croisant les codes et les répondants \cite{miles2003} ne produit de
résultat que lue en lignes, sur ce qui converge, ce qui sépare et ce qui n'apparaît qu'une
fois. Un constat est dit \textbf{saturé} lorsque les cinq récits le documentent,
\textbf{clivant} lorsque le matériau se scinde entre deux postures, \textbf{isolé} lorsqu'un
seul répondant le rapporte. Il est alors conservé, mais sans portée générale. Le mot
«~saturé~» vaut dans son acception faible~: sur cinq entretiens codés après la clôture de la
collecte, la saturation théorique \cite{glaser1967} n'est pas revendiquée.

Le suivi du nombre de codes nouveaux apportés par chaque entretien, dans l'ordre de
passation, précise cette réserve. Le premier entretien fait apparaître vingt-sept codes,
le deuxième onze, le troisième quatre, le quatrième deux et le cinquième quatre~;
quarante-huit des cinquante-trois codes de la grille sont mobilisés au total. La
décroissance est nette et la structure de la grille se stabilise dès le troisième
entretien, les cinq axes étant tous mobilisés dès le premier récit. Aucun entretien
n'atteint toutefois zéro, et le dernier apporte encore quatre codes inédits, dont deux relèvent de l'axe~C~: la saturation n'est pas atteinte, et la moitié de ce qu'il
apporte encore relève de l'axe que la section~4.4 déclare le plus
fragile.

Le tableau~\ref{tab:statut-resultats} récapitule le statut des principaux résultats au regard du panel ; les sous-sections qui suivent les développent axe par axe.

\begin{table}[h]
\centering
\small
\begin{tabular}{|p{0.46\linewidth}|p{0.15\linewidth}|p{0.25\linewidth}|}
\hline
\textbf{Résultat} & \textbf{Codes} & \textbf{Statut} \\
\hline
Zone sanctuarisée, et convergence de sa frontière & A2.1, A2.2 & Saturé \\
\hline
Seuil au-delà duquel l'outil coûte plus qu'il ne rapporte & A4.1, A4.3 & Saturé \\
\hline
Temps gagné réinvesti dans la vérification & A4.1, B2.3 & Saturé, une exception \\
\hline
Apprentissage de l'outil autonome ou entre pairs & A3.1, A3.2 & Saturé, une exception \\
\hline
Déclaration ou non de l'usage au client & C2.2, C2.1 & Saturé, une exception \\
\hline
Valeur déplacée vers le jugement et la responsabilité & C3.1, C3.2 & Saturé \\
\hline
Absence de contrôle dédié aux contenus générés & D2.3 & Saturé \\
\hline
Cadre normatif jugé insuffisant & D1.3, D5.2 & Saturé \\
\hline
Effet perçu sur l'apprentissage du junior & B2.2, T1.2 & Clivant \\
\hline
Adaptation de la pratique de revue & D2.2, T3.2 & Clivant \\
\hline
Interpellation du cabinet par un client & C2.3 & Isolé \\
\hline
Transparence assumée vis-à-vis du client & C2.1 & Isolé \\
\hline
\end{tabular}
\caption{Statut des principaux résultats au regard du panel}
\label{tab:statut-resultats}
\end{table}

\subsection{Axe A --- L'appropriation des outils}

Les tâches déléguées sont les mêmes d'un bout à l'autre du panel et relèvent toutes du
traitement d'un matériau déjà constitué~: reformulation, synthèse, structuration d'un plan,
traduction, premier jet sur un sujet peu maîtrisé. Aucun ne décrit une délégation de
l'analyse elle-même~; la formule d'E3 résume la position commune, «~assistant, pas
rédacteur~».

La sanctuarisation est saturée, et sa frontière l'est aussi. La zone que chacun ne délègue
pas se définit par la sensibilité des données et par la nécessité d'un contexte que l'outil
n'a pas connu, non par la difficulté de la tâche. Ce sont les données chiffrées
confidentielles pour J1 («~dealbreaker~»), l'analyse conduite après un exercice de crise
pour J2, au motif noté que «~l'IA n'était pas dans la salle~». Pour E3, ce sont
l'architecture de sécurité, les vulnérabilités et les données classifiées réservées à la
«~main humaine~»~; pour E2, les sujets spécifiques et sensibles~; pour E1, les choix
engageant une recommandation ou une orientation. Les missions diffèrent~; la frontière,
elle, ne varie pas.

Le seuil de spécificité est le deuxième constat saturé, et il est chiffré par tous. J1
compte vingt minutes de reformulations de consignes sans résultat exploitable puis quinze
minutes de rédaction directe, sur un sujet où l'outil renvoyait à des textes inexistants.
J2 compte dix minutes puis cinq. Pour E3, quarante-cinq minutes de reprise sur une
reformulation où deux concepts clés avaient été intervertis, contre vingt minutes en
production directe. Pour E1, une nuit à reprendre un livrable client intégralement produit
avec assistance. E2 énonce le même seuil par son complément, estimant à «~80~\% du travail~» la part de gain net
et désignant comme compétence propre le fait de «~savoir où pertinent, où pas~».

Le troisième constat saturé contredit l'intuition ordinaire du gain de productivité~: le
temps gagné ne sort pas de la boucle de production, il se déplace de la rédaction vers la
vérification. Le déplacement est chiffré~: deux heures en trois relectures pour J1, deux
heures en recoupement croisé des sources pour E3, une heure en contrôle de cohérence pour J2,
E1 énonçant que le gain devrait y être affecté. L'exception vient du répondant le plus
utilisateur~: E2 réinvestit ce temps dans la relation client, le c\oe{}ur du métier étant
selon ses notes de comprendre et non de produire.

L'apprentissage de l'outil se fait enfin par essai personnel ou entre pairs~: essais
personnels pour J1, apprentissage sur le tas pour J2, autodidaxie pour E2, carnet de
consignes constitué par tâtonnement pour E3. Un dispositif interne existe pourtant, et il
est obligatoire~: la formation d'environ une heure décrite au chapitre~1 conditionne
l'ouverture de l'accès à l'outil homologué. Deux répondants l'évoquent~: E1, sous la forme
d'une session de sensibilisation qu'il n'a pas approfondie, et J1, sous celle d'une session
rappelant les règles d'usage de l'outil homologué et d'un webinaire qu'une mission l'a
empêché de suivre. Aucun ne la désigne comme le canal par lequel il a appris à se servir de
l'outil. Le corpus confirme ainsi par son silence ce que le chapitre~1 établit par la
description du dispositif~: l'encadrement porte sur l'entrée dans l'usage, non sur l'usage
lui-même. Une exception doit cependant être signalée, et elle vient de l'extérieur du
cabinet. E3 suit depuis plusieurs mois une certification externe, engagée de sa propre
initiative et motivée par la demande des clients. Il note qu'elle vaut autant comme document
à valoriser que comme apprentissage, et il y revient plus tard dans l'entretien, en évoquant
sa formation continue.
L'observation a une portée qui excède l'axe~: la compétence sur ces outils commence à se
certifier hors du cabinet, ce qui rejoint le déplacement de la valeur discuté en 4.2.

\subsection{Axe B --- Apprentissage et transmission}

Une règle traverse les trois encadrants, énoncée depuis trois anciennetés différentes~:
l'exécution manuelle initiale n'est pas négociable, l'accélération vient après. E1 estime que
les juniors doivent avoir fait les tâches ingrates une ou deux fois à la main, faute de quoi
ils ne sauront pas juger ce que l'outil produit. E2 décrit la même séquence et ajoute que
déléguer sans avoir jamais produit interdit d'évaluer. E3 déplace le critère sur la
définition du métier, un bon consultant étant celui qui sait ce qu'il met dans ses livrables
et non celui qui sait formuler une consigne.

L'effet sur le raisonnement du junior est en revanche clivant. J1 documente sur lui-même ce
que la grille prévoyait d'observer. Il se dit «~plus relecteur que rédacteur~» et qualifie
d'«~illusion de compétence~» la maîtrise apparente que donne un support produit vite. Il
formule son inquiétude sous forme d'échéance, se demandant s'il saura encore, dans six mois,
«~faire structure seul~». J2, à ancienneté comparable, ne partage pas ce diagnostic. La
différence tient vraisemblablement à la nature des tâches, créatives et non factuelles chez
lui. 
Côté encadrant, le constat est convergent~: E1 décrit des «~livrables très propres très vite~»
suivis d'une interrogation sur la profondeur, et des juniors «~très bons forme, mais mal
justifier fond~»~; E3 parle d'une illusion de maîtrise sans avoir creusé.

Le résultat le plus net de cet axe est négatif. L'appauvrissement de l'interaction entre
junior et encadrant, inscrit en B3.2 à titre d'hypothèse, n'est documenté par personne, et
les deux positions décrivent le mouvement inverse. J1 observe que les commentaires de
l'encadrant portent désormais sur la solidité du raisonnement plutôt que sur l'alignement des
titres, et relève lui-même le caractère paradoxal de cette amélioration de la relecture
senior. E1 décrit la relecture comme transmission du jugement plutôt que correction
technique, affirme que cette transmission devient plus importante avec l'outil, et range la
formation des juniors à la vérification critique parmi les composantes importantes de son
rôle. Ce résultat est repris en 4.2.

\subsection{Axe C --- Légitimité et relation client}

La question de la déclaration de l'usage au client est documentée par les cinq récits, et
quatre y répondent par la non-déclaration, avec la même
justification~: le client achète un résultat, non un processus. J1 parle de sujet tabou~; J2
juge la mention sans apport. E3 la juge non pas dissimulatrice mais contre-productive, le
client payant une expertise et non un outil. E1, interrogé par des clients, répond que
l'analyse finale et les recommandations reviennent aux équipes, ce qui est exact. Mais il
indique qu'il ne dirait jamais qu'un support a été produit avec l'outil. E2 est isolé, et cet
isolement est cohérent avec sa fonction~: il aborde le sujet de sa propre initiative en
rendez-vous commercial et en fait un argument de positionnement.

L'interpellation par le client est le point le plus mince du corpus. Un seul répondant, E1,
rapporte une situation où un client a soulevé la question, et une seule conséquence
commerciale d'une défaillance est documentée~: des erreurs factuelles non détectées, une
méfiance non exprimée directement, une mission suivante sur laquelle un autre consultant a
été demandé. La question projective sur l'horizon de dix ans produit en revanche une
convergence complète, les cinq plaçant la valeur future du cabinet du côté du jugement et de
la responsabilité assumée. Hormis le récit d'E1, tout cet axe relève du déclaratif projectif.

\subsection{Axe D --- Gouvernance, contrôle et risques}

Ce que les répondants connaissent est un cadre de confidentialité et une liste d'outils
homologués, non une doctrine d'usage. Le chapitre~1 rappelle qu'il n'existe pas de politique
écrite en la matière. Tous le jugent insuffisant, pour deux raisons opposées. Pour les uns,
la règle est floue sur son périmètre~: J1 la décrit comme un flou total s'agissant des
outils non homologués et qualifie ce flou de «~dangereux~», E1 juge les règles
«~un peu légères~». Pour les autres, elle est claire mais muette sur l'essentiel. E2, qui la
trouve «~trop timide~», observe qu'elle autorise et interdit des outils sans rien dire de la
manière de vérifier ce qu'ils produisent. E3 la trouve générale au point d'appliquer une
règle plus stricte. Un cadre
existe donc~: il délimite des outils, il ne délimite pas des pratiques.

L'absence de contrôle dédié suit. Le circuit de revue hiérarchique fonctionne et personne ne
le conteste, mais aucun des cinq ne décrit d'étape propre aux contenus produits avec
assistance. J1 est le plus explicite~: personne ne demande d'où vient un contenu, le contrôle
porte sur le fond et non sur le processus, et les encadrants «~soupçonnent mais ne posent
pas~». E3 confirme qu'aucune procédure spécifique n'existe et rappelle que la charge repose
sur celui qui signe.

Les incidents, eux, existent~: un livrable client entier comportant chiffres faux et
références inventées, repris en une nuit, pour E1~; plusieurs occurrences pour E3~; un
quasi-incident pour J1, sur un document individuel où des projets inexistants avaient été
introduits. Leur traitement constitue le résultat de gouvernance le plus important du corpus.
D'une part, les incidents ne remontent pas. Celui de J1 a été raconté à un collègue de même
niveau, sur le registre de la plaisanterie, et n'a jamais été porté à la connaissance de son
encadrante. Les notes consignent le motif~: la honte et la crainte d'un reproche
d'inattention. D'autre part, le contrôle s'adapte, mais individuellement et après coup. E1
contrôle désormais le processus et non plus le seul contenu final, demandant ce qui relève du
collaborateur, ce qui relève de l'outil et si les sources ont été vérifiées. E3 a formalisé
pour lui-même un contrôle en trois temps~: cohérence interne, contrôle croisé des sources,
test de défendabilité devant le client. E2, qui n'a pas vécu d'incident marquant, indique ne
pas contrôler différemment un contenu assisté. Aucune de ces adaptations n'est écrite,
partagée ni opposable.

\subsection{Les deux clivages}

Le premier a été rencontré axe par axe~: les juniors et les encadrants observent le même
phénomène et ne se le disent pas. J1 s'inquiète de son propre apprentissage, E1 de la
profondeur des livrables qu'il reçoit, et le corpus ne comporte aucune trace d'une
conversation entre ces deux positions. Le silence est lui-même documenté, dans des termes
presque identiques et sans concertation~: J1 note que tout le monde utilise l'outil et que
personne ne le dit, et rattache ce silence à la crainte du reproche de ne plus réfléchir. J2
emploie l'expression de «~secret de polichinelle~». Ce non-dit interne, codé en T3.1, n'était
pas anticipé. La grille prévoyait l'opacité vis-à-vis du client, non la dissimulation
entre collègues. Le résultat n'est pas un désaccord entre positions, c'est l'absence de
l'échange qui permettrait d'en avoir un.

Le second est plus inattendu~: les encadrants ne forment pas un bloc. E1, manager de dix ans,
n'utilise presque pas l'outil, se déclare plus rapide seul et emploie le terme de «~béquille~»
à propos des juniors. E2, associé de vingt-cinq ans et sponsor de la practice
concernée, l'utilise quotidiennement et le décrit comme un «~accélérateur de compétence, pas
substitut~». La posture ne suit donc ni l'ancienneté ni le grade, et l'écart se lit dans la
pratique de contrôle décrite ci-dessus, ce qu'enregistre le code T3.2.

Un dernier élément doit être signalé sans être tranché. Le chapitre~1 décrit un cadre
reposant sur un seul outil homologué. Or E2 précise ne jamais verser de données
confidentielles dans des outils non officiels, formulation qui suppose un recours à ces outils
pour les contenus qui ne sont pas confidentiels. J1, lui, déclare un usage officieux dont il doute que
la prudence qui l'accompagne soit générale. La règle n'est donc pas comprise de la même façon
selon la position occupée. Statuer sur la conformité de ces pratiques dépasse ce que le
dispositif permet d'établir. Le constat est autre~: un cadre diversement interprété
n'encadre pas.

\section{Confrontation aux cadres théoriques}

\subsection{Le seuil de spécificité, ou la frontière dentelée décrite de l'intérieur}

Dell'Acqua et ses coauteurs \cite{dellacqua2026} établissent, sur une population de
consultants, l'existence d'une frontière irrégulière séparant les tâches sur lesquelles
l'outil produit un gain net de celles sur lesquelles il dégrade la performance, et montrent
qu'elle est difficile à percevoir de l'intérieur. Le corpus décrit exactement cette
frontière, et la décrit cinq fois, par des personnes qui n'ont pas lu l'article.

Deux nuances font l'intérêt de la confrontation. Les répondants ont localisé la frontière, là
où l'expérimentation la donne pour peu visible. Mais la méthode est unique et coûteuse~:
l'échec chiffré en minutes perdues puis mémorisé en règle personnelle. La localisation est
acquise~; elle a été payée cinq fois, séparément, par des personnes qui ignoraient que leurs
collègues avaient payé le même prix. Ensuite, cette frontière ne sépare pas les tâches faciles
des tâches difficiles, mais celles dont le résultat se vérifie par recoupement de celles qui
exigent un contexte non observé. C'est la frontière de la sanctuarisation décrite en 4.1~:
les répondants n'ont donc pas construit deux règles mais une seule.

Un désalignement mérite enfin d'être noté. Là où Noy et Zhang \cite{noy2023} puis
\mbox{Brynjolfsson}, Li et Raymond \cite{brynjolfsson2025} montrent des gains concentrés sur les
moins expérimentés, le gain le plus élevé est ici déclaré par le répondant le plus ancien. La
variable discriminante est le portefeuille de tâches plutôt que le grade.

\subsection{Un contrôle qui s'adapte par l'incident et non par la règle}

Kellogg, Valentine et Christin \cite{kellogg2020} analysent le contrôle algorithmique comme un
terrain de lutte, en recensant les leviers par lesquels une technologie oriente et discipline
le travail et les résistances qu'ils suscitent. Le terrain observé se laisse nommer par défaut
dans cette grille~: il n'y a pas ici de contrôle algorithmique du travail, mais un travail
assisté sur lequel aucun dispositif de contrôle n'a été construit. Le cadre existant est un
cadre d'autorisation (tel outil est permis, tel autre non) et non un cadre de
vérification. Le contournement se produit donc en amont de toute règle contraignante~: le
non-dit codé en T3.1 est social et non procédural, puisqu'aucune procédure ne prescrit de
déclarer quoi que ce soit.

Lebovitz, Lifshitz-Assaf et Levina \cite{lebovitz2022} éclairent le versant individuel, en
montrant comment des professionnels arbitrent face à l'opacité d'un outil au moment de porter
un jugement qui les engage. Le protocole en trois temps d'E3 est exactement un dispositif de
ce type. La règle qu'il transmet aux juniors en est la maxime~:
«~si tu peux pas justifier, mets pas~». Ce dispositif existe donc, il est fin et efficace~;
mais il est privé,
construit par un individu à ses frais, et rien dans l'organisation ne le rend disponible aux
autres. La difficulté que soulèvent Lebovitz, Levina et Lifshitz-Assaf \cite{lebovitz2021} à
propos de l'étalon de vérification se retrouve sous forme pratique~: ce qui a sauvé J1 lors de
son quasi-incident est d'avoir conservé une source de vérité à côté de lui, circonstance
heureuse et non dispositif.

Le mécanisme d'ensemble peut alors être énoncé. Le circuit de revue du cabinet est hérité de
la standardisation des qualifications propre à la bureaucratie professionnelle
\cite{mintzberg1982}, où le contrôle porte sur la personne formée plutôt que sur l'acte~; il
tient tant que le relecteur sait quoi chercher. Or la nature des erreurs a changé (E1
relève que celles de l'outil sont plus subtiles que les erreurs humaines) sans que le
contenu de la qualification ait suivi. La compétence de vérification s'acquiert donc par
l'expérience, et cette expérience est celle de l'incident.

\subsection{La relation junior-encadrant ne s'appauvrit pas~: elle change d'objet}

Le résultat le plus solide du corpus est celui qui contredit l'hypothèse formulée en amont. La
grille avait inscrit en B3.2 l'appauvrissement de l'interaction entre junior et encadrant,
cohérent avec ce qu'établit Anthony \cite{anthony2021} sur des analystes juniors qui, faute
d'interroger le fonctionnement d'une technologie opaque, développent une expertise plus
superficielle. Le mécanisme attendu était direct~: si la machine produit le premier jet, la
boucle de correction par le senior perd son objet et le compagnonnage s'étiole.

Ce n'est pas ce que le terrain donne à voir, et les deux positions le contredisent
indépendamment. J1 a pourtant toutes les raisons de noircir le tableau, puisqu'il est
celui qui parle d'«~illusion de compétence~». Il décrit sa première proposition commerciale
comme un moment de formation fait de commentaires de titres et de structure, et observe que
ces commentaires portent aujourd'hui sur la solidité du raisonnement. E1 décrit
symétriquement ce que transmet une relecture (le jugement, la raison pour laquelle un
chiffre est suspect) et l'estime plus importante encore avec l'outil.

L'interprétation la plus économe passe par le modèle de conversion des connaissances
\cite{nonaka1997}. L'hypothèse d'appauvrissement supposait que la socialisation, mode de
conversion du tacite vers le tacite, avait pour support le premier jet du junior~: la machine
produisant ce jet, le support disparaissait. Le corpus indique qu'il n'a pas disparu mais
changé de niveau. La forme étant déléguée, la revue se reporte sur ce qu'elle n'atteignait
auparavant qu'après plusieurs itérations~: la justification, la pertinence, la source. La
socialisation ne se réduit pas~; elle porte sur un objet de rang supérieur. Et ce qu'E1 décrit
(énoncer ce qui rend un chiffre suspect, ce qui doit être vérifié et pourquoi) relève
moins de la socialisation que de l'externalisation, conversion du tacite vers l'explicite~:
l'outil oblige l'encadrant à mettre en mots un jugement jusque-là incorporé dans sa pratique.
Le gain de transmissibilité est réel, et il est involontaire.

La portée du résultat doit être bornée. La revue doit avoir lieu et porter sur le processus,
ce qu'E1 fait depuis son incident et qu'E2 déclare ne pas faire~: le bénéfice tient à la
pratique de l'encadrant plutôt qu'à la technologie, et il n'est pas généralisé dans le panel.
Il coexiste ensuite avec une perte en amont que le même E1 signale, le chemin de réflexion
initial n'étant plus parcouru. Cette perte affecte l'internalisation, conversion de
l'explicite vers le tacite qui suppose l'exécution. La relation monte d'un cran~; ce sur quoi elle
s'exerce a été appauvri en amont. Le corpus documente enfin des perceptions et non des
acquisitions.

C'est ici que la règle des trois encadrants prend son sens théorique. Exécuter manuellement
une fois avant d'accélérer est la condition d'internalisation. Sans exécution, le savoir
explicite déposé dans les méthodologies ne se convertit pas en savoir tacite. Le
professionnel se retrouve alors dans la situation, nommée en B4.2, de valider un contenu
qu'il n'aurait pas su produire. Or la standardisation des qualifications, mécanisme de
coordination principal d'une bureaucratie professionnelle \cite{mintzberg1982}, est dans le
conseil produite par le cabinet lui-même faute d'institution externe qui la délivre
\cite{nordenflycht2010}. Ce que la section 2.1.2 établissait sur le plan structurel se
vérifie donc ici~: la perturbation du parcours d'apprentissage n'atteint pas une fonction
support, elle atteint le mécanisme de coordination lui-même.

\subsection{Un point plus fragile~: la légitimité}

La confrontation à Abbott \cite{abbott1988} et à Suchman \cite{suchman1995} appelle plus de
réserve, faute de matériau. Les cinq répondants situent spontanément la valeur future du
cabinet là où Abbott place le noyau défendable d'une juridiction, dans l'inférence plutôt que
dans la collecte ou la mise en forme. Ceux qui refusent de mentionner l'outil au client
protègent un signal plutôt qu'une information, ce qui rejoint la lecture d'Alvesson
\cite{alvesson2001}. Mais cette convergence a été obtenue sur une question projective à dix
ans, dont la forme appelle une réponse valorisante. Les registres de Suchman permettent de
classer ce qui a été dit. Le registre cognitif est documenté par les cinq, chacun rapportant
la confiance du client à la fiabilité et à la maîtrise du domaine. Le moral l'est par trois,
par la responsabilité assumée devant le client, à quoi s'ajoute la non-déclaration de
l'usage, qui porte sur ce qui a été réellement fait et par qui. Le pragmatique n'est
documenté qu'une fois.
Ils ne permettent pas, en revanche, d'établir comment la légitimité se reconstruit devant
un client.

\section{Recommandations managériales}

Les recommandations qui suivent ne procèdent pas d'une opinion sur l'opportunité de ces
outils~: elles reprennent ce que le terrain a formulé, en le rendant opérationnel. Interrogés
séparément sur la décision qu'ils prendraient à la place de la direction, les cinq répondants
ont convergé sans s'être parlé et sur le même point d'équilibre. Aucun ne demande
d'interdiction, tous demandent que la vérification soit organisée, et E1 et E3 y ajoutent la
traçabilité des contenus produits avec assistance. Sur un
sujet où le clivage entre ceux qui produisent et ceux qui contrôlent était attendu, cette
convergence est elle-même un résultat.

\subsection{Séquencer l'autorisation d'usage dans le parcours du junior}

Chaque type de mission identifie une liste courte de tâches fondatrices (plan de livrable,
synthèse d'un corpus documentaire, chiffrage, note de cadrage) que le collaborateur produit
une fois sans assistance, la production étant relue et validée comme telle. Le passage est
acté~; l'usage de l'outil est ensuite non seulement autorisé mais encouragé sur ces mêmes
tâches, conformément à la séquence décrite par E2. La mesure doit se présenter comme une
séquence et non comme une restriction, faute de quoi elle produira les contournements
qu'anticipe la littérature sur les interventions organisationnelles \cite{venkatesh2008}.
Une règle vécue comme bridante sur une population déjà équipée est contournée en silence, et
le cabinet perd à la fois la formation et l'information.

\subsection{Rendre la vérification déclarative et graduer la relecture}

La mesure proposée par E1 est directement transposable~: le livrable produit avec assistance
porte une mention déclenchant une relecture renforcée. Elle répond à un constat saturé,
l'absence de toute étape de contrôle propre aux contenus générés. Trois conditions en
déterminent le succès. Le dispositif doit être déclaratif et non détectif~: il sert à
permettre le signalement de l'usage, non à le rechercher. Il doit être non punitif et énoncé
comme tel, faute de quoi il renforcera le non-dit interne au lieu de le lever. Le silence
actuel est motivé par la crainte du reproche, et un marquage perçu comme un aveu produirait
l'effet inverse. Il doit enfin être associé à la question qu'E1 pose déjà en revue, qui
distingue ce qui relève du collaborateur de ce qui relève de l'outil et s'assortit de la
vérification des sources. La finalité, telle qu'E3 la formule, est de responsabiliser et non
de contrôler.

\subsection{Constituer la vérification en compétence formée et outillée}

Quatre des cinq demandent une formation, et tous précisent qu'elle ne doit pas porter sur la
manipulation de l'outil. J1 la veut sur la détection des erreurs, le risque n'étant pas que
l'outil se trompe mais que le consultant ne le voie pas. E1 veut que les encadrants soient
formés à une relecture spécifique, les erreurs de l'outil étant plus subtiles. E2 la range du
côté de la compétence et non de la défiance~; E3 la veut obligatoire. Le cabinet n'a pas à en
concevoir le contenu, lequel existe déjà sous forme privée dans le corpus. Le protocole en
trois temps d'E3 est un standard éprouvé et transmissible en une page, que complète la
pratique de J1 consistant à conserver une source de vérité indépendante à côté du document en
production. La mesure consiste à expliciter ce savoir tacite~: rassembler les protocoles
individuels, les formaliser, les enseigner. La charte d'une page demandée par J2, qui
distingue ce qui est possible, ce qui ne l'est pas et ce qui doit être vérifié, en constitue
le format~; le canal d'échange que J2 et E2 réclament l'un et l'autre en constitue le
support d'entretien.

\subsection{Faire évoluer le dispositif autrement que par l'accident}

Le principe directeur qui commande les mesures précédentes s'énonce simplement~: un dispositif
de contrôle qui n'évolue que par accident est un dispositif qui attend l'accident. Le corpus
l'établit deux fois~: les encadrants qui ont modifié leur revue l'ont fait après un incident
vécu, et le seul quasi-incident rapporté par un junior n'est jamais remonté. Deux mesures en
découlent. La première est un traitement organisé des quasi-incidents, dissocié de toute
conséquence individuelle. Un incident détecté avant le client est une information de valeur,
et sa remontée doit être plus avantageuse que son silence. Les cas recueillis alimentent la
révision périodique du protocole de vérification, de sorte que l'apprentissage cesse de
dépendre de la répartition aléatoire des incidents. La seconde est la clarification du
périmètre normatif, dont le flou est signalé par les deux positions et produit un usage dont
l'ampleur échappe au cabinet. Cette évolution est enfin cohérente avec le repositionnement que
les cinq décrivent déjà (la valeur vendue se déplace vers le jugement et la responsabilité
assumée \cite{suchman1995}), car un dispositif de vérification explicite est ce qui rend
cette promesse vérifiable.

\section{Limites des résultats}

Les limites du dispositif ont été énoncées au chapitre précédent~; celles qui suivent portent
sur les résultats eux-mêmes.

La première tient à l'assise empirique, décrite en 3.4 et non reprise ici. Sa conséquence
sur les résultats est en revanche précise~: les convergences relevées ne valent pas comme
fréquences. Un constat dit saturé signifie que les cinq récits le documentent, non qu'il
est répandu dans la practice. Un constat dit clivant oppose des répondants nommément
identifiés, non deux groupes constitués, et le clivage sur l'adaptation de la revue passe à
l'intérieur du groupe des encadrants.

La deuxième tient à l'inégale solidité des axes. L'axe C est le plus mince et porte pourtant
la troisième question directrice. Un seul répondant rapporte une interpellation réelle par un
client, un seul une conséquence commerciale documentée. Le reste est du déclaratif projectif,
recueilli sur des questions dont la forme appelle une réponse valorisante. La convergence
obtenue sur cet axe est donc d'une autre nature que celle obtenue sur les axes A et D, où elle
repose sur des récits d'incidents datés et chiffrés. Elle dit ce que le métier devrait être,
non comment la légitimité se reconstruit devant un client.

La troisième tient au matériau. Les extraits reproduits ici sont des fragments de notes prises
en séance, donnés tels qu'ils ont été consignés ou rapportés au style indirect~: aucun n'a
valeur de citation littérale, et l'affectation d'un code dépend d'une formulation en partie
celle de l'enquêteur.

La quatrième tient à la position de l'enquêteur, décrite en 3.4, et affecte spécifiquement les deux clivages.
Cette position a probablement rendu possible le recueil
du non-dit codé en T3.1~: un enquêteur extérieur n'aurait guère obtenu de deux juniors
l'aveu d'un usage dissimulé à leur hiérarchie. Mais elle expose au risque symétrique de
reconnaître d'autant plus aisément un phénomène qu'on s'attendait à trouver. Le silence entre
positions est établi par ce que les répondants en disent, non par l'observation de leurs
échanges.

La cinquième porte sur le cadre normatif~: le dispositif ne permet ni d'établir l'ampleur des
usages relevés en 4.1, ni de statuer sur leur conformité. Ce point est versé au dossier comme
une observation de gouvernance appelant vérification, non comme un constat de manquement.

% ==============================================================================
% Fin du chapitre 4
% ==============================================================================